% !TEX root = ../../../main.tex
% Block card: detection.spike_v1 -- "Spike detection (v1)" (Part I, the blocks chapter).
% This card and its entry are the worked example AUTHORING.md points to.
\cardfield{What it does}
Marks each moment where the signal stands clearly out from its own local average, up or
down, and counts a tight cluster of such moments as one spike, placed at its most extreme
point. Each spike it keeps is one detection.

\cardfield{When to use it}
When you want a quick, explainable list of fast spikes on one channel, found by the rule
\citet{adamatzky2023multiscalar} used on \emph{Schizophyllum commune}, with a few changes
that make it work on recordings whose baseline sits far from zero. It is a good first pass
before measuring the trains and gaps between spikes with \blockref{interrogation.trains} or
\blockref{interrogation.intervals}.

\cardfield{What goes in, what comes out}
In goes a span of one channel. Out comes one short stretch per spike, centred on it, with
the spike's own moment marked inside. The chain row counts them and says how many went up
and how many went down.

\cardfield{Settings that matter}
Tune \emph{delta\_mv} first: how far, in millivolts, a sample must stand from its local
average to count. Raise it if the block fires on noise, lower it if it misses spikes you can
see. The block's page tells you how \emph{delta\_mv} compares with this span's own
wobble, and turns amber when it is under four times that. Then \emph{width}: how wide the
local average is. Changes slower than that window are invisible to the block, so use
\blockref{detection.drop_detection} for slow drops. Then \emph{d}: candidates closer
together than \emph{d} samples are one spike. \emph{polarity} keeps rises, dips or both.

\cardfield{Watch out for}
The default \emph{delta\_mv} of 1\,mV is a placeholder, chosen for one recording's noise
level; on a noisier channel it fires on noise. The block converts millivolts by the
recording's declared unit; if no unit is declared it assumes volts and says so on the chain
row, and on a recording that is really in millivolts that makes the threshold a thousand
times too small. \emph{width} and \emph{d} count samples, not seconds. Several spikes closer
together than \emph{d} come out as one.
